
\Formula{A$_{\bm2}$}{
	计算
	\Equation{
	    \lim_{x\to 0}\frac{\rme^{-1/x^{2}}}{x}\ .
	}
}
\Proof{证明I\score{1.5}}{
	\par 我们有
	\Equation{
		0\leqslant\frac{\rme^{-1/x^{2}}}{x}\leqslant\frac{\braI{1+\sfrac{1}{x^{2}}}^{-1}}{x}=\frac{x}{x^{2}+1}\ ,
	}
	使用夹逼准则即得
	\Equation{
		\lim_{x\to 0}\frac{\rme^{-1/x^{2}}}{x}=\boxed{0}\ .
	}
}
\Proof{证明II\score{1.5}}{
	\par 注意到
	\Equation{
		\frac{\rmd}{\rmd x}\braI{\frac{\rme^{-1/x^{2}}}{x}}=-\frac{\rme^{-1/x^{2}}\braI{x^{2}-2}}{x^{4}}\ ,
	}
	因此该函数在 $\braI{-\sqrt{2}\and 0}$ 和 $\braI{0\and\sqrt{2}}$ 上递增且分别有上界和下界 $0$ , 于是由单调有界准则知其极限存在.
	\par 设
	\Equation{
		\lim_{x\to 0}\frac{\rme^{-1/x^{2}}}{x}=A\ ,
	}
	则作代换 $t=-x$ 又有
	\Equation{
		\lim_{t\to 0}-\frac{\rme^{-1/t^{2}}}{t}=A\ ,
	}
	于是 $A=-A$ , 即 $A=\boxed{0}$ .
}
\Proof{证明III\score{1}}{
	\Equation{
		\lim_{x\to 0^{+}}\frac{\rme^{-1/x^{2}}}{x}=\lim_{x\to 0^{+}}\exp\braI{-\frac{x^{2}\ln x+1}{x^{2}}}\xlongequal[\text*{(指数)}]{\mathrm{L'H}}\lim_{x\to 0^{+}}\exp\braI{-\frac{2\ln x+1}{2}}=0\ ,\\[-14mm]
	}
	\Equation{
		\lim_{x\to 0^{-}}\frac{\rme^{-1/x^{2}}}{x}=-\lim_{x\to 0^{+}}\frac{\rme^{-1/x^{2}}}{x}=0\ ,
	}

\newpage

	\noindent 于是
	\Equation{
		\lim_{x\to 0}\frac{\rme^{-1/x^{2}}}{x}=\boxed{0}\ ,
	}
	其中 $\lim_{x\to 0^{+}}x^{2}\ln x=-1$ 是因为
	\Equation{
		\lim_{x\to 0^{+}}x^{2}\ln x=\lim_{x\to 0^{+}}\frac{\ln x}{\sfrac{1}{x^{2}}}\xlongequal{\mathrm{L'H}}\lim_{x\to 0^{+}}-\frac{x^{2}}{2}=0\ .
	}
}
\Proof{证明IV\score{1}}{
	\Equation{
		\lim_{x\to 0}\frac{\rme^{-1/x^{2}}}{x}=\lim_{x\to\infty}\frac{x}{\rme^{x^{2}}}\xlongequal{\mathrm{L'H}}\lim_{x\to\infty}\frac{1}{2x\rme^{x^{2}}}=\boxed{0}\ .
	}
}
\Proof{要点}{
	+ 如果无法直接对表达式使用L'Hôpital法则, 可以尝试作倒数代换.
	+ 使用L'Hôpital法则时要注意分子分母是否符合要求, 例如对证明III中指数部分的处理.
}

\bigskip

\Formula{B$_{\bm2}$}{
	计算
	\Equation{
	    \lim_{x\to 0}\frac{\ln\braI{1+x}+\ln\braI{1-x}-\ln\cos^{x}x}{x^{2}}\ .
	}
}
\Proof{证明I\score{1.5}}{
	\Equation*{
		\lim_{x\to 0}\frac{\ln\braI{1+x}+\ln\braI{1-x}-\ln\cos^{x}x}{x^{2}}&=\lim_{x\to 0}\frac{\ln\braI{1-x^{2}}}{x^{2}}-\lim_{x\to 0}\frac{\ln\cos x}{x}\\[1mm]
		&=-\lim_{x\to 0^{+}}\frac{\ln\braI{1+x}-\ln 1}{x}-\lim_{x\to 0}\frac{\ln\cos x-\ln\cos 0}{x}\\[1mm]
		&=-\frac{\rmd}{\rmd t}\braI{\ln t}\subst{t=1}{}-\frac{\rmd}{\rmd t}\braI{\ln\cos t}\subst{t=0}{}\\[-1mm]
		&=\boxed{-1}\ .
	}
}
\Proof{证明II\score{1}}{
	\Equation{
		\lim_{x\to 0}\frac{\ln\braI{1+x}+\ln\braI{1-x}-\ln\cos^{x}x}{x^{2}}&=\lim_{x\to 0}\frac{\ln\braI{1-x^{2}}}{x^{2}}-\lim_{x\to 0}\frac{\ln\cos x}{x}\\[1mm]
		&=\lim_{x\to 0}\frac{-x^{2}}{x^{2}}-\lim_{x\to 0}\frac{\cos x-1}{x}\\[-1mm]
		&=\boxed{-1}\ .
	}
}

\newpage

\Proof{要点}{
	+ 注意将 $x^{2}$ 代换为 $x$ 后 $x\to 0$ 转变为 $x\to 0^{+}$ .
}

\bigskip

\Formula{C$_{\bm2}$}{
	计算
	\Equation{
	    \lim_{n\to\infty}\braI{n-\frac{1}{n}\sum_{k=1}^{n}\sqrt{n^{2}+k}}\ .
	}
}
\Proof{证明I\score{2}}{
	\par 我们有
	\Equation{
		n-\frac{1}{n}\sum_{k=1}^{n}\sqrt{n^{2}+k}=-\frac{1}{n}\sum_{k=1}^{n}\frac{k}{n+\sqrt{n^{2}+k}}\geqslant-\sum_{k=1}^{n}\frac{k}{2n^{2}}=-\frac{n+1}{4n}\ ,\\[-14mm]
	}
	\Equation{
		n-\frac{1}{n}\sum_{k=1}^{n}\sqrt{n^{2}+k}=-\frac{1}{n}\sum_{k=1}^{n}\frac{k}{n+\sqrt{n^{2}+k}}\leqslant-\sum_{k=1}^{n}\frac{k}{2n^{2}+n}=-\frac{n+1}{4n+2}\ ,
	}
	使用夹逼准则即得
	\Equation{
		\lim_{n\to\infty}\braI{n-\frac{1}{n}\sum_{k=1}^{n}\sqrt{n^{2}+k}}=\boxed{-\frac{1}{4}}\ .
	}
}
\Proof{证明II\score{2}}{
	\Equation{
		\lim_{n\to\infty}\braI{n-\frac{1}{n}\sum_{k=1}^{n}\sqrt{n^{2}+k}}&=\lim_{n\to\infty}\frac{1}{n}\sum_{k=1}^{n}\braI{n-\sqrt{n^{2}+k}}\\[1mm]
		&=-\lim_{n\to\infty}\frac{1}{n}\sum_{k=1}^{n}\frac{k}{n+\sqrt{n^{2}+k}}\\[1mm]
		&=\lim_{n\to\infty}\frac{1}{n}\sum_{k=1}^{n}\braI{\frac{k}{2n}-\frac{k}{n+\sqrt{n^{2}+k}}}-\lim_{n\to\infty}\frac{1}{n}\sum_{k=1}^{n}\frac{k}{2n}\\[1mm]
		&=\lim_{n\to\infty}\sum_{k=1}^{n}\frac{k^{2}}{2n^{2}\braI{n+\sqrt{n^{2}+k}}^{2}}-\lim_{n\to\infty}\frac{n+1}{4n}\\[-1mm]
		&=\boxed{-\sfrac{1}{4}}\ ,
	}
	其中
	\Equation{
		\sum_{k=1}^{n}\frac{k^{2}}{2n^{2}\braI{n+\sqrt{n^{2}+k}}^{2}}\leqslant\sum_{k=1}^{n}\frac{1}{8n^{2}}=\frac{1}{8n}\implies\lim_{n\to\infty}\sum_{k=1}^{n}\frac{k^{2}}{2n^{2}\braI{n+\sqrt{n^{2}+k}}^{2}}=0\ .
	}
}

\newpage

\Proof{证明III\score{1.5}}{
	\par 分别考虑 $n=1$ 和 $n=2$ 时的\link{Alpha}{Lagrange余项}, 我们有
	\Equation{
		\sqrt{1+\frac{k}{n^{2}}}=1+\frac{k}{2n^{2}}-\frac{k^{2}}{8n^{4}}\braI{1+\xi_{1}}^{-3/2}\ ,\quad\xi_{1}\in\braI{0\and\frac{k}{n^{2}}}\\[-14mm]
	}
	\Equation{
		\sqrt{1+\frac{k}{n^{2}}}=1+\frac{k}{2n^{2}}-\frac{k^{2}}{8n^{4}}+\frac{3k^{3}}{48n^{6}}\braI{1+\xi_{2}}^{-5/2}\ ,\quad\xi_{2}\in\braI{0\and\frac{k}{n^{2}}}
	}
	于是
	\Equation{
		-\sum_{k=1}^{n}\frac{k}{2n^{2}}\leqslant\sum_{k=1}^{n}\braI{1-\sqrt{1+\frac{k}{n^{2}}}}\leqslant-\sum_{k=1}^{n}\braI{\frac{k}{2n^{2}}-\frac{k^{2}}{8n^{4}}}\ ,\\[-14mm]
	}
	\Equation{
		-\frac{n+1}{4n}\leqslant\sum_{k=1}^{n}\braI{1-\sqrt{1+\frac{k}{n^{2}}}}\leqslant-\frac{12n^{3}+10n^{2}-3n-1}{48n^{3}}\ ,
	}
	使用夹逼准则即得
	\Equation{
		\lim_{n\to\infty}\braI{n-\frac{1}{n}\sum_{k=1}^{n}\sqrt{n^{2}+k}}=\lim_{n\to\infty}\sum_{k=1}^{n}\braI{1-\sqrt{1+\frac{k}{n^{2}}}}=\boxed{-\frac{1}{4}}\ .
	}
}
\Proof{证明IV\score{1.5}}{
	\par 考虑\link{Alpha}{Peano余项}, 我们有
	\Equation{
		\sqrt{1+\frac{k}{n^{2}}}=1+\frac{k}{2n^{2}}+o\braI{n^{-1}}\ ,
	}
	于是
	\Equation{
		\lim_{n\to\infty}\braI{n-\frac{1}{n}\sum_{k=1}^{n}\sqrt{n^{2}+k}}&=\lim_{n\to\infty}\sum_{k=1}^{n}\braI{1-\sqrt{1+\frac{k}{n^{2}}}}\\[1mm]
		&=\lim_{n\to\infty}\sum_{k=1}^{n}\braI{-\frac{k}{2n^{2}}+o\braI{n^{-1}}\!}\\[1mm]
		&=-\lim_{n\to\infty}\braI{\frac{n+1}{4n}+o\braI{1}\!}\\[-1mm]
		&=\boxed{-\sfrac{1}{4}}\ .
	}
}
\Proof{要点}{
	+ 证明II体现了重要的拟合思想: 将复杂的表达式拟合为简单的形式, 然后将问题转化为验证该拟合在极限意义下无损失, 即两种表达式的差是无穷小量.
}

\newpage
